%! Setting Up a Test for the Difference of Two Population Means — Unit 7, Lesson 7.8 — Lesson Plan
\documentclass[10pt]{article}
\usepackage{apstats-article}
\usepackage{apstats-boxes}
\usepackage{graphicx}
\graphicspath{{images/}}
\newcommand{\TallMath}[1]{$\displaystyle #1\rule[-1.4em]{0pt}{3.2em}$}
\newcommand{\UnitNumberName}{Unit 7: Inference for Quantitative Data: Means \quad}
\newcommand{\LessonNumberName}{Lesson 7.8: Setting Up a Test for the Difference of Two Population Means}

\begin{document}
\begin{center}
    {\Large\bfseries \CourseName: \SchoolYear \\
    {\small\bfseries \UnitNumberName \LessonNumberName}}
\end{center}

\begin{tcolorbox}[colback=sky, colframe=navy, boxrule=0.9pt, arc=2mm,
    left=2mm, right=2mm, top=1.5mm, bottom=1.5mm]
    \textbf{Primary Objective:} Students will identify when a two-sample $t$-test for a
    difference of two population means is appropriate, write null and alternative hypotheses
    in terms of $\mu_1 - \mu_2$, and verify the conditions for inference --- extending the
    one-sample $t$-test framework to two independent groups. (Big Idea: VAR)
\end{tcolorbox}

\renewcommand{\arraystretch}{1.6}
\begin{skillbox}[Priority Ideas \& Skills]{goldbox}
\begin{minipage}[t]{0.48\linewidth}
    \textbf{AP Skill 1 --- Selecting Statistical Methods}
    \begin{itemize}
        \item Identify a two-sample $t$-test as the appropriate method for testing
              a claim about the difference of two population means when $\sigma_1$
              and $\sigma_2$ are unknown (Skill 1.E; VAR-7.F).
        \item Identify null and alternative hypotheses for a two-sample $t$-test in
              terms of $\mu_1 - \mu_2$ (Skill 1.F; VAR-7.G).
    \end{itemize}
    \textbf{AP Skill 4 --- Statistical Argumentation}
    \begin{itemize}
        \item Verify that conditions for inference apply: independence for both
              samples and approximately normal sampling distribution of
              $\bar x_1 - \bar x_2$ (Skill 4.C; VAR-7.H).
    \end{itemize}
\end{minipage}
\hfill
\begin{minipage}[t]{0.48\linewidth}
    \textbf{Key Understandings}
    \begin{itemize}
        \item When comparing two independent groups and both $\sigma$'s are unknown,
              use a \textbf{two-sample $t$-test for $\mu_1 - \mu_2$}.
        \item $H_0: \mu_1 - \mu_2 = 0$ (equivalently $\mu_1 = \mu_2$); the alternative
              reflects the direction of the research question.
        \item Two sets of conditions to verify: (1) each sample is random or from a
              randomized experiment; if sampling without replacement, apply the 10\%
              rule to \emph{each} sample; (2) the sampling distribution of
              $\bar x_1 - \bar x_2$ is approximately normal --- check both samples
              separately.
        \item Distinguish two-sample $t$-test (independent groups) from matched-pairs
              $t$-test (paired data on the same subjects).
    \end{itemize}
\end{minipage}
\end{skillbox}

\begin{skillbox}[{Vocabulary, Concepts, \& Theorems}]{greenbox}
\begin{minipage}[t]{0.48\linewidth}
    \begin{itemize}
        \item \textbf{Two-sample $t$-test for $\mu_1 - \mu_2$} --- a significance test
              comparing two independent groups when both population standard deviations
              are unknown
        \item \textbf{Null hypothesis ($H_0$)} --- $H_0: \mu_1 - \mu_2 = 0$
              (equivalently $H_0: \mu_1 = \mu_2$)
        \item \textbf{Alternative hypothesis ($H_a$)} --- $H_a: \mu_1 - \mu_2 < 0$,
              $> 0$, or $\ne 0$; the direction is determined by the research question,
              set \emph{before} seeing data
        \item \textbf{Independent samples} --- data from two unrelated groups; no
              natural pairing between observations
    \end{itemize}
\end{minipage}
\hfill
\begin{minipage}[t]{0.48\linewidth}
    \begin{itemize}
        \item \textbf{10\% condition (two samples)} --- if sampling without replacement,
              $n_1 \le 10\%N_1$ \emph{and} $n_2 \le 10\%N_2$
        \item \textbf{Normal condition (two samples)} --- if both $n_1 \ge 30$ and
              $n_2 \ge 30$, CLT applies; if either is $< 30$, check \emph{both}
              samples for strong skewness and outliers
        \item \textbf{Two-sided vs.\ one-sided alternative} --- two-sided ($\ne$) tests
              whether the means differ; one-sided ($<$ or $>$) tests a directional claim
        \item \textbf{Two-sample vs.\ matched pairs} --- independent groups $\to$
              two-sample $t$; paired data $\to$ one-sample $t$ on differences
    \end{itemize}
\end{minipage}
\end{skillbox}

\begin{skillbox}[Activate Prior Knowledge \& Spiral Review]{sky}
\begin{minipage}[t]{0.55\linewidth}\vspace{0pt}
    \textbf{Spiral Review (warm-up)}
    \begin{itemize}
        \item One-sample $t$-test setup: identifying the method, writing $H_0/H_a$,
              verifying conditions --- from Lesson 7.4
        \item Two-sample $t$-interval conditions: each sample random, 10\% rule for
              each group, Normal condition for both --- from Lesson 7.6
        \item Writing hypotheses in terms of population parameters, not statistics ---
              connecting Lessons 7.4 and 7.6
    \end{itemize}
\end{minipage}
\hfill
\begin{minipage}[t]{0.38\linewidth}\vspace{0pt}\centering
    % Authored warm-up compiles to target/ — no source PDF to embed here.
\end{minipage}
\end{skillbox}

\begin{skillbox}[Hook]{sky}
\begin{minipage}[t]{0.55\linewidth}
    \textbf{Which Training Program Works Better? (3--4 min)}\\
    A hospital randomly assigns patients recovering from knee surgery to one of two
    physical therapy programs: standard PT or an accelerated PT protocol. After 8 weeks,
    researchers record the number of days until each patient returns to full mobility.\\[4pt]
    Ask: ``How would you determine whether one program leads to faster recovery
    \emph{on average}?''\\[4pt]
    Key discussion points:
    \begin{itemize}
        \item We have two independent groups --- not paired data.
        \item We want to compare $\mu_1$ (standard PT) vs.\ $\mu_2$ (accelerated PT).
        \item Before computing anything: choose the method, write hypotheses, verify conditions.
    \end{itemize}
\end{minipage}
\hfill
\begin{minipage}[t]{0.40\linewidth}
    \textbf{Bridge from Lesson 7.6}\\
    In Lesson 7.6 we built a \emph{confidence interval} for $\mu_1 - \mu_2$:\\[4pt]
    Two-sample $t$-interval: estimate the plausible range of $\mu_1 - \mu_2$.\\[6pt]
    In Lesson 7.8 we run a \emph{significance test} for $\mu_1 - \mu_2$:\\[4pt]
    Two-sample $t$-test: ``Is there convincing evidence that $\mu_1 \ne \mu_2$?''\\[6pt]
    Same conditions; different purpose.
\end{minipage}
\end{skillbox}

\begin{skillbox}[Lesson: Choosing the Method]{sky}
\begin{minipage}[t]{0.48\linewidth}
    \textbf{When to Use a Two-Sample $t$-Test}
    \begin{enumerate}
        \item Quantitative response variable measured on \emph{two independent groups}.
        \item Population standard deviations $\sigma_1$ and $\sigma_2$ are unknown
              (virtually always true in practice).
        \item Goal: test a claim about $\mu_1 - \mu_2$.
    \end{enumerate}
    \textit{PT example: days to full mobility is quantitative; standard PT and accelerated PT
    are two independent groups; $\sigma$'s unknown $\Rightarrow$ two-sample $t$-test.}
\end{minipage}
\hfill
\begin{minipage}[t]{0.48\linewidth}
    \textbf{Two-Sample vs.\ Matched-Pairs --- Key Distinction}
    \begin{itemize}
        \item \textbf{Matched pairs:} each observation in group 1 is naturally paired
              with one in group 2 (same subject before/after, or two measurements on
              the same unit). $\to$ one-sample $t$-test on differences.
        \item \textbf{Two independent samples:} different subjects randomly assigned
              to or selected from two separate groups. $\to$ two-sample $t$-test.
    \end{itemize}
    \textit{Ask: ``Is there a natural pairing?'' If yes $\to$ matched pairs. If no $\to$
    two-sample.}
\end{minipage}
\end{skillbox}

\begin{skillbox}[Lesson (cont.): Hypotheses \& Conditions]{sky}
\begin{minipage}[t]{0.48\linewidth}
    \textbf{Writing Hypotheses (VAR-7.G)}
    \begin{itemize}
        \item Always in terms of \emph{population} parameters $\mu_1$ and $\mu_2$,
              never sample statistics.
        \item $H_0: \mu_1 - \mu_2 = 0$ \quad (equivalently $\mu_1 = \mu_2$)
        \item Choose $H_a$ based on the research question --- \emph{before} data:
              \begin{itemize}
                  \item ``Is Group 1 \emph{lower} on average?''
                        $\to H_a: \mu_1 - \mu_2 < 0$
                  \item ``Is Group 1 \emph{higher} on average?''
                        $\to H_a: \mu_1 - \mu_2 > 0$
                  \item ``Do the groups \emph{differ}?''
                        $\to H_a: \mu_1 - \mu_2 \ne 0$
              \end{itemize}
    \end{itemize}
    \textit{PT example (let $\mu_1$ = standard, $\mu_2$ = accelerated):
    $H_0: \mu_1 - \mu_2 = 0$; $H_a: \mu_1 - \mu_2 > 0$ (standard takes longer).}
\end{minipage}
\hfill
\begin{minipage}[t]{0.48\linewidth}
    \textbf{Verifying Conditions (VAR-7.H)}
    \begin{enumerate}
        \item \textbf{Random/Independence (both samples):}
              \begin{itemize}
                  \item Each group: random sample or randomized experiment.
                  \item If sampling without replacement: $n_1 \le 10\%N_1$
                        \emph{and} $n_2 \le 10\%N_2$.
              \end{itemize}
        \item \textbf{Normal (sampling dist.\ of $\bar x_1 - \bar x_2$):}
              \begin{itemize}
                  \item Both $n_1 \ge 30$ and $n_2 \ge 30$: CLT applies to each.
                  \item If \emph{either} $n < 30$: check \emph{both} sample
                        distributions for strong skewness and outliers.
              \end{itemize}
    \end{enumerate}
    \textit{PT: $n_1 = 18$, $n_2 = 17$ --- both $< 30$, so examine plots of
    \emph{both} groups.}
\end{minipage}
\end{skillbox}

\begin{skillbox}[Active Monitoring]{redbox}
\begin{minipage}[t]{0.48\linewidth}
    \textbf{Circulate and Check}
    \begin{itemize}
        \item Students write $\mu_1 - \mu_2$ (not $\bar x_1 - \bar x_2$) in hypotheses.
        \item $H_0$ is $\mu_1 - \mu_2 = 0$; null value is always zero in standard setup.
        \item Alternative direction matches the research question; chosen before data.
        \item Normal condition check applies to \emph{both} samples independently.
        \item Students correctly distinguish independent samples from matched pairs.
    \end{itemize}
\end{minipage}
\hfill
\begin{minipage}[t]{0.48\linewidth}
    \textbf{Cold-Call Prompts}
    \begin{itemize}
        \item ``Why is the null value for $\mu_1 - \mu_2$ always 0?''
        \item ``What would change in the setup if the same patients received
              both treatments?''
        \item ``If $n_1 = 40$ but $n_2 = 18$, which Normal condition check applies?''
        \item ``A student writes $H_0: \bar x_1 = \bar x_2$. What's wrong?''
    \end{itemize}
\end{minipage}
\end{skillbox}

\begin{skillbox}[Group Work \& Differentiation]{redbox}
\begin{multicols}{3}
    \textbf{Tier R --- Remediate}
    \begin{itemize}
        \item Provide a side-by-side template: ``$H_0: \mu_1 - \mu_2 = 0$;
              $H_a: \mu_1 - \mu_2$ (circle: $<$ / $>$ / $\ne$) $0$'' and the
              two-condition checklist.
        \item Focus on identifying independent vs.\ paired designs.
        \item Activity: match scenarios to hypotheses only (Tier R items).
    \end{itemize}
    \columnbreak
    \textbf{Tier A --- Approaching}
    \begin{itemize}
        \item Write full setup (method, hypotheses, conditions) for each scenario.
        \item Justify each condition with numbers from the context.
    \end{itemize}
    \columnbreak
    \textbf{Tier E --- Extension}
    \begin{itemize}
        \item Complete all scenarios, then design an original two-sample study.
        \item Explain the trade-offs of a one-sided vs.\ two-sided alternative for
              the same research context.
    \end{itemize}
\end{multicols}
\end{skillbox}

\begin{skillbox}[Individual Work \& Assessment]{redbox}
\begin{minipage}[t]{0.48\linewidth}
    \textbf{Exit Ticket (5 min)}\\
    A data scientist believes the mean commute time for workers using public transit differs
    from the mean commute time for workers who drive. She randomly surveys 22 transit users
    and 25 drivers. Dotplots of both samples show roughly symmetric distributions with no
    outliers.
    \begin{enumerate}
        \item Name the appropriate inference method.
        \item State $H_0$ and $H_a$ in terms of $\mu_1 - \mu_2$.
        \item Verify conditions for both samples.
    \end{enumerate}
\end{minipage}
\hfill
\begin{minipage}[t]{0.48\linewidth}
    \textbf{AP-Style Check}\\
    \textit{A researcher randomly assigns 30 volunteers to Diet A and 30 to Diet B. After
    12 weeks he records weight lost (kg). Both samples are approximately normal. He tests
    $H_0: \mu_A - \mu_B = 0$ vs.\ $H_a: \mu_A - \mu_B \ne 0$. Which method is correct?}
    \begin{enumerate}[label=(\Alph*)]
        \item One-sample $t$-test; the data come from one experiment.
        \item Matched-pairs $t$-test; both groups share the same time period.
        \item Two-sample $t$-test for $\mu_A - \mu_B$. \textcolor{keyred}{\textbf{$\leftarrow$ correct}}
        \item $z$-test; both sample sizes equal 30.
    \end{enumerate}
\end{minipage}
\end{skillbox}

\begin{skillbox}[Reinforcement \& Extension]{goldbox}
\begin{minipage}[t]{0.48\linewidth}
    \textbf{Homework Overview}
    \begin{itemize}
        \item Four to five contexts across varied settings: identify the method,
              state hypotheses for $\mu_1 - \mu_2$, and verify conditions for both
              samples. No $t$-statistic computation yet.
        \item Students write a one-sentence justification for the direction of
              $H_a$ in each problem.
    \end{itemize}
\end{minipage}
\hfill
\begin{minipage}[t]{0.48\linewidth}
    \textbf{Extension / Preview}
    \begin{itemize}
        \item \textit{Extension:} A student chooses $H_a: \mu_1 - \mu_2 > 0$
              after seeing $\bar x_1 > \bar x_2$ in the data. Explain what is
              wrong and why this inflates the Type I error rate.
        \item \textit{Preview:} Lesson 7.9 carries out the two-sample $t$-test ---
              compute the $t$-statistic, find the $p$-value using technology, and
              write a conclusion in context.
    \end{itemize}
\end{minipage}
\end{skillbox}

\end{document}
